\section{Introduction}
\label{sec:introduction}

\subsection{Purpose}
This document is the Software Requirements Specification (SRS) and accompanying
technical report for \textbf{Campus ERP}, a production-oriented, multi-tenant
Enterprise Resource Planning system built for the SEN4121 (Large System
Environment) project brief. It specifies the functional and non-functional
requirements of the system, documents the architecture that satisfies the
brief's mandatory microservices constraints, and reports the verification
evidence collected against those requirements. The document is structured to
follow ISO/IEC/IEEE~29148 guidance for content, adapted to the length and
audience of a course deliverable.

\subsection{Scope}
Campus ERP is decomposed into four independently deployable microservices --
\textbf{Identity}, \textbf{Academic}, \textbf{Marketing \& Finance}, and
\textbf{Administration \& Human Resources} -- each with its own PostgreSQL
database, fronted by a single NGINX API gateway, and connected for one
cross-module workflow via asynchronous messaging over RabbitMQ. The system
supports multiple independent institutions (tenants) from one deployment,
with every tenant's data fully isolated at the database layer. The scope
covers the three domains required by the brief (Academic; Marketing \&
Finance; Administration \& HR), the mandatory security, data, messaging and
DevOps constraints, and the supporting React single-page application.

Explicitly out of scope for this submission: real mobile-money settlement
(a signed, disclosed simulator stands in for MTN MoMo / Orange Money),
real biometric hardware (a signed, time-boxed QR challenge is used instead,
as permitted by the brief), multi-factor authentication and Kubernetes
manifests (both listed as optional bonus work), and any legally binding
certification that the implemented payroll rates satisfy current
Cameroonian tax law -- see the statutory rate governance note in
\cref{sec:hr-module}.

\subsection{Definitions, Acronyms and Abbreviations}
\begin{longtable}{L{3.3cm}L{11cm}}
  \toprule
  \rowcolor{ceTeal}\thd{Term} & \thd{Definition} \\
  \midrule
  \endhead
  Tenant & An institution using Campus ERP; all tenant-owned data is isolated
    by \texttt{tenant\_id} and PostgreSQL row-level security. \\
  RLS & Row-Level Security -- a PostgreSQL feature that restricts which rows a
    query may see or modify based on the active session's tenant context. \\
  RBAC & Role-Based Access Control. \\
  JWT & JSON Web Token, used for short-lived API access tokens. \\
  Outbox / Inbox & A reliability pattern where a producer writes an event in
    the same database transaction as its business change (outbox), and a
    consumer records processed event IDs to guarantee idempotent handling
    (inbox). \\
  CNPS & Caisse Nationale de Pr\'evoyance Sociale -- Cameroon's national social
    insurance fund (pension, family allowance, work-injury contributions). \\
  IRPP / PAYE & Imp\^ot sur le Revenu des Personnes Physiques -- Cameroon's
    personal income tax, withheld monthly (Pay-As-You-Earn). \\
  CAC & Centimes Additionnels Communaux -- an additional council tax levied as
    a percentage of IRPP. \\
  FCFA / XAF & Franc CFA (BEAC), the currency used throughout the system. \\
  WAT & West Africa Time (UTC+1), the timezone used for all displayed dates. \\
  PITR & Point-In-Time Recovery. \\
  RPO / RTO & Recovery Point / Time Objective. \\
  \bottomrule
\end{longtable}

\subsection{References}
\begin{itemize}[leftmargin=1.4em]
  \item ISO/IEC/IEEE 29148:2018, \textit{Systems and software engineering --
    Life cycle processes -- Requirements engineering}.
  \item SEN4121 Large System Environment, Summer 2026 project brief (course
    handout), Faculty of Information and Communication Technologies.
  \item OWASP Top 10:2021, \textit{The Ten Most Critical Web Application
    Security Risks}.
  \item Direction G\'en\'erale des Imp\^ots (DGI), \textit{Code G\'en\'eral des
    Imp\^ots} and 2026 Finance Law circular -- see Appendix~\ref{app:sources}
    for the exact publications consulted for payroll rates.
  \item PostgreSQL 17, RabbitMQ 3.13, FastAPI, React and Vite official
    documentation (implementation references).
\end{itemize}

\subsection{Document Overview}
\cref{sec:overall-description} gives the product perspective, user classes and
constraints. \cref{sec:architecture} presents the microservices architecture.
\cref{sec:security} and \cref{sec:data-design} cover security and data design.
\cref{sec:uml-suite} contains the full UML diagram suite. Sections~7--9 walk
through the Academic, Finance \& Marketing, and Administration \& HR modules
with functional requirements and screenshots. \cref{sec:nfr-testing-devops}
reports non-functional, testing and DevOps evidence. \cref{sec:limitations}
states honest limitations and the academic-integrity disclosure. The
appendices hold the requirement traceability matrix and cited sources.
